\chapter{Quando il nemico è in casa: i bias più pericolosi}

Non tutti i bias sono uguali. Alcuni sono fastidiosi ma relativamente innocui (come sopravvalutare le nostre capacità di guida), altri possono avere conseguenze devastanti per individui e società.

Il bias di conferma, per esempio, non è solo la tendenza a cercare informazioni che confermano le nostre idee: è il meccanismo che crea le "camere dell'eco" dove prosperano teorie complottiste, estremismi politici e pseudoscienza. In un mondo dove ciascuno può scegliere le proprie fonti di informazione, questo bias ancestrale sta letteralmente frammentando la realtà condivisa.

L'avversione alle perdite ci fa aggrappare a investimenti fallimentari, relazioni tossiche e lavori che ci rendono infelici. L'effetto gregge ci spinge a seguire mode finanziarie che finiscono sempre in bolle speculative.

Ma forse il bias più pericoloso di tutti è quello che ci fa credere di essere immuni dai bias. Pensare "io sono troppo intelligente per cascarci" è spesso il primo passo per cascarci in pieno.

\bigskip

Se dovessimo stilare una classifica dei bias più devastanti dell'era moderna, il bias di conferma si aggiudicherebbe senza dubbio il primo posto. Ma per capire perché è diventato così letale, dobbiamo prima capire come ha cambiato volto nell'era digitale.

Nei capitoli precedenti abbiamo visto come il bias di conferma fosse un meccanismo di sopravvivenza perfettamente sensato: se i vostri antenati avevano imparato che certi frutti erano velenosi, era meglio non sperimentare ogni volta. Il bias di conferma garantiva coerenza e stabilità alle convinzioni che avevano dimostrato il loro valore pratico.

Ma c'era un limite naturale a questo meccanismo: il mondo fisico non mente. Se continuavate a cercare cibo dove non ce n'era, morivate di fame. Se continuavate a bere da una fonte avvelenata, morivate avvelenati. La realtà fisica funzionava come un sistema di correzione automatica che impediva al bias di conferma di diventare completamente delirante.

Nel mondo dell'informazione digitale, questo sistema di correzione è saltato. Potete cercare conferme alle vostre convinzioni all'infinito senza mai scontrarvi con una realtà che vi contraddice. Anzi, gli algoritmi di raccomandazione dei social media e dei motori di ricerca sono progettati proprio per darvi sempre più di quello che già credete.

Il risultato è che persone intelligenti, istruite, razionali in tutti gli altri aspetti della vita, possono sviluppare convinzioni completamente scollegate dalla realtà e mantenerle per anni senza mai incontrare informazioni che le mettano in discussione.

\medskip

Prendiamo l'esempio delle teorie complottiste. Una volta, se avevate una teoria bizzarra su come funzionava il mondo, eravate limitati alle persone del vostro paese, della vostra cerchia sociale, della vostra biblioteca locale. Prima o poi, qualcuno vi avrebbe detto: "Ma questa cosa non ha senso".

Oggi, invece, qualsiasi teoria, per quanto assurda, può trovare online una comunità di migliaia di persone che la pensano allo stesso modo. Il bias di conferma, combinato con la vastità di Internet, crea quello che possiamo chiamare "delirium collettivi": gruppi di persone che si confermano a vicenda convinzioni sempre più estreme e distaccate dalla realtà.

È il bias di conferma che ha preso steroidi. Non si limita più a cercare informazioni che confermano le vostre idee: crea interi ecosistemi informativi alternativi dove quelle idee non vengono mai messe in discussione.

E una volta che siete dentro uno di questi ecosistemi, uscirne diventa incredibilmente difficile. Ogni informazione che contraddice le vostre convinzioni viene automaticamente interpretata come prova del complotto, propaganda del nemico, fake news. Il bias di conferma non si limita a selezionare le informazioni: le reinterpreta attivamente per renderle coerenti con quello che già credete.

\bigskip

Ma il bias di conferma è solo il più spettacolare di una serie di bias che, nell'era moderna, sono diventati molto più pericolosi di quanto non fossero nell'ambiente ancestrale. Prendiamo l'avversione alle perdite: la tendenza a odiare perdere qualcosa molto più di quanto amiamo guadagnare la stessa cosa.

Nell'ambiente ancestrale, questo bias aveva perfettamente senso. Perdere risorse poteva significare morte per fame o per freddo. Era letteralmente più importante non perdere quello che si aveva che guadagnare qualcosa in più. Ma nel mondo moderno, l'avversione alle perdite spesso ci impedisce di cogliere opportunità che potrebbero migliorare drasticamente la nostra vita.

Quante persone rimangono bloccate in lavori che detestano perché hanno paura di perdere la sicurezza del posto fisso? Quante restano in relazioni infelici perché il dolore di perdere quello che hanno sembra maggiore della possibilità di trovare qualcosa di meglio? Quanti investitori tengono azioni in perdita per anni, sperando che risalgano, invece di vendere e investire in qualcosa di più promettente?

L'avversione alle perdite ci trasforma in prigionieri del nostro status quo, anche quando quello status quo ci rende infelici. È come se il nostro cervello fosse programmato per preferire una certezza mediocre a un'incertezza potenzialmente fantastica.

\medskip

E poi c'è l'effetto gregge, forse il bias più sottovalutato ma potenzialmente più distruttivo di tutti. Nell'ambiente ancestrale, seguire il gruppo era quasi sempre la strategia più sicura. Se tutti scappavano in una direzione, probabilmente c'era un buon motivo. Se tutti evitavano un certo cibo, probabilmente era velenoso.

Ma nel mondo moderno, "tutti lo stanno facendo" può significare "tutti stanno per cadere nella stessa trappola". L'effetto gregge è il motore delle bolle finanziarie, delle mode dannose, dei panici morali collettivi.

Pensate alla bolla delle dot-com alla fine degli anni '90. Persone intelligenti, esperti di finanza, investitori professionali versavano milioni in aziende che non avevano mai fatto un centesimo di profitto, solo perché "tutti" stavano investendo nelle tecnologie Internet. Il meccanismo mentale era esattamente lo stesso che spingeva i nostri antenati a seguire la migrazione della tribù: se tutti vanno in quella direzione, deve esserci un buon motivo.

Il problema è che nel mondo moderno, "tutti" possono sbagliare contemporaneamente su scale che i nostri antenati non avrebbero mai potuto immaginare. Quando una tribù di 50 persone seguiva una direzione sbagliata, si perdevano 50 persone. Quando una società di milioni di persone segue una direzione sbagliata, le conseguenze possono essere catastrofiche.

\bigskip

Ma c'è un bias ancora più insidioso di tutti quelli che abbiamo visto finora: il bias del punto cieco. È il bias che ci impedisce di vedere i nostri bias. È la convinzione che, mentre tutti gli altri sono influenzati da pregiudizi e distorsioni cognitive, noi siamo diversi. Noi siamo obiettivi, razionali, immuni da queste debolezze umane.

Questo bias è particolarmente forte nelle persone più istruite e intelligenti. Più sappiamo sui bias cognitivi, più tendiamo a pensare di esserne immuni. È come se la conoscenza teorica dei nostri limiti cognitivi ci desse un falso senso di sicurezza che ci rende paradossalmente più vulnerabili.

È l'effetto Dunning-Kruger applicato ai bias cognitivi stessi: un po' di conoscenza ci rende più confidenti, ma non necessariamente più bravi. Anzi, la confidenza eccessiva può renderci meno attenti proprio ai segnali che dovrebbero metterci in guardia.

Quante volte avete sentito qualcuno dire: "Io non sono influenzabile dalla pubblicità", proprio mentre indossava vestiti di marca, guidava un'auto scelta per il suo valore simbolico, e aveva in tasca uno smartphone scelto più per status che per funzionalità?

Il bias del punto cieco ci fa credere che i nostri pregiudizi siano intuizioni, che le nostre distorsioni siano saggezza, che i nostri errori sistematici siano eccezioni casuali. È l'ultimo meccanismo di difesa dell'ego: anche quando dobbiamo ammettere che gli umani in generale sono irrazionali, riusciamo sempre a convincerci di essere l'eccezione.

\medskip

Ma forse l'aspetto più pericoloso di tutti questi bias è come si rinforzano a vicenda, creando quello che possiamo chiamare "tempeste cognitive perfette". Il bias di conferma ci fa cercare informazioni che confermano le nostre idee. L'effetto gregge ci spinge a circondarci di persone che la pensano come noi. L'avversione alle perdite ci impedisce di abbandonare convinzioni in cui abbiamo già investito emotivamente. E il bias del punto cieco ci convince che tutto questo sia perfettamente razionale.

È come avere un sistema immunitario che, invece di proteggere il corpo dalle infezioni, protegge le nostre convinzioni dalla realtà. Una volta che questo sistema si attiva, diventa incredibilmente difficile da fermare.

Prendiamo l'esempio della polarizzazione politica moderna. Due persone che partono con opinioni leggermente diverse su un argomento possono finire, nel giro di pochi anni, a vivere in universi informativi completamente separati, dove ciascuna è convinta che l'altra sia vittima di propaganda e disinformazione.

Il bias di conferma le spinge a cercare fonti che confermano le loro opinioni. L'effetto gregge le porta a frequentare persone che rinforzano quelle opinioni. L'avversione alle perdite le rende riluttanti ad ammettere di aver sbagliato. E il bias del punto cieco le convince che sono loro quelle razionali, mentre gli altri sono stati "lavati del cervello".

Il risultato finale è la frammentazione della società in tribù cognitive che non riescono più a comunicare tra loro, ciascuna convinta di essere l'unica a vedere la realtà chiaramente.

\bigskip

Cosa possiamo fare di fronte a questi "bias di massa distruzione"? La prima tentazione è sempre quella di cercare di eliminarli, di diventare perfettamente razionali, di superare i nostri limiti cognitivi attraverso la forza di volontà e l'educazione.

Ma questa strategia è destinata al fallimento, perché si basa su un fraintendimento fondamentale di cosa siano i bias cognitivi. Non sono difetti che possono essere corretti: sono caratteristiche così fondamentali del nostro modo di pensare che eliminarle significherebbe smettere di essere umani.

L'approccio più realistico è sviluppare quello che possiamo chiamare "umiltà cognitiva": la consapevolezza che i nostri processi di pensiero sono sistematicamente imperfetti, e che questa imperfezione ci rende vulnerabili a errori prevedibili.

Non si tratta di diventare scettici paranoici che dubitano di tutto. Si tratta di riconoscere che i nostri istinti cognitivi, per quanto utili siano stati nell'ambiente ancestrale, possono tradirci nel mondo moderno. E di sviluppare strategie per verificare le nostre convinzioni, diversificare le nostre fonti di informazione, e esporci deliberatamente a punti di vista che mettono in discussione le nostre certezze.

\medskip

Forse la lezione più importante che possiamo imparare dai bias più pericolosi è che la razionalità non è uno stato naturale che dobbiamo raggiungere, ma una disciplina che dobbiamo praticare. È come la forma fisica: non è qualcosa che si ottiene una volta per tutte, ma qualcosa che va mantenuta con esercizio costante.

E come per la forma fisica, l'esercizio deve essere mirato ai nostri punti deboli. Se siete particolarmente suscettibili al bias di conferma, dovete deliberatamente cercare fonti che contraddicono le vostre opinioni. Se tendete a seguire il gregge, dovete allenarvi a essere sospettosi quando "tutti" fanno la stessa cosa. Se soffrite di avversione alle perdite, dovete praticare il lasciare andare cose che non vi servono più.

Non sarà mai naturale, perché va contro milioni di anni di programmazione evolutiva. Ma è possibile, ed è necessario se vogliamo navigare un mondo che è cambiato molto più velocemente del nostro cervello.

Alla fine, riconoscere che abbiamo il nemico in casa - che i nostri peggiori nemici sono spesso i nostri stessi processi di pensiero - non è motivo di disperazione. È motivo di speranza. Perché una volta che sappiamo chi è il nemico, possiamo iniziare a difenderci.

Non vinceremo mai completamente questa battaglia. Ma possiamo almeno smettere di perdere tutte le guerre per distrazione.